\section{Background and Related Work} 
\label{sec:bg}

\subsection{LLM Serving Systems} 
\label{subsec:bg-llm}

% \jiahuan{add a state diagram: waiting, running}

% \jiahuan{TODO(P1): need to explain what is ``iteration'' in the background section}

Numerous systems have been proposed to improve LLM serving efficiency, including advanced batching strategies for throughput optimization~\cite{fang2021turbotransformers,yu2022orca}, memory management techniques such as PagedAttention~\cite{kwon2023efficient}, CPU offloading~\cite{xu2024pie,sheng2023flexgen}, and vAttention~\cite{prabhu2025vattention}, and GPU kernel-level optimizations for accelerating attention and decoding~\cite{aminabadi2022deepspeed, dao2023flashattention2,fastertransformer,wang2021lightseq,shah2024flashattention}. To better utilize available resources and improve scheduling, model parallelism and pipelining frameworks have been introduced~\cite{aminabadi2022deepspeed,li2023alpaserve,pope2023efficiently}, along with parameter sharing mechanisms to reduce memory and computation overhead~\cite{sheng2023s}. Speculative decoding has emerged as a promising technique to reduce tail latency by predicting likely tokens ahead of time and verifying them in parallel~\cite{miao2024specinfer,stern2018blockwise}. Additionally, systems like FastServe~\cite{wu2023fast} explore preemptive scheduling policies to reduce job completion time (JCT) in multi-tenant serving environments. While these systems reduce energy consumption as a byproduct of latency or throughput improvements, \aryan{we treat energy as a first-class optimization target and specifically exploit the control opportunities exposed by P/D-disaggregated serving.} \name directly targets energy efficiency through frequency and routing-aware scheduling, which has been relatively underexplored despite its critical implications for sustainable deployment. 

\subsection{Energy-Efficient LLM Serving} 
\label{sec:bg-energy}

Several works have started to explore energy-efficient and power management frameworks for LLM inference~\cite{stojkovic2025dynamollm,stojkovic2025tapas,li2025ecoserve,reddy2025ai,fernandez2025energy,stojkovic2024towards}.
For example, DynamoLLM~\cite{stojkovic2025dynamollm} shows benefits of GPU frequency control in LLM inference serving based on request characteristics (predicted input/output tokens) to yield energy savings.
It proposes a hierarchical design to efficiently route incoming requests to respective GPU pools that vary by model sharding and frequency while changing frequency on a coarse-grained basis of load characteristics.
ThrottLL'eM~\cite{kakolyris2025throttll} predicts future KV cache usage and batch size to reduce frequencies and instance sizes for energy efficiency.
Similarly, $\mu$-Serve~\cite{qiu2024power} shows benefits of dynamic frequency scaling and proposes a model-serving framework to optimize for power consumption by co-serving multiple ML models. \aryan{In contrast, our approach is explicitly phase-aware: it distinguishes prefill and decode as separate control regimes with different latency sensitivities and energy behaviors, rather than applying a unified serving-time power policy. Moreover, \name also performs per-engine-iteration frequency selection within the inference engine to react to short-term workload fluctuations under latency SLOs.}

\aryan{Taken together, these works show that DVFS and scheduling can improve energy and system efficiency, but they do not treat P/D disaggregation as an energy-control surface, nor do they study the phase-specific energy–frequency behavior and decode-side batch-boundary effects that motivate \name.}
\sout{
These works show early signals towards power-saving opportunities by coarse-grained control in LLM inference but do not explore the significant benefits of \cameraready{per-engine-iteration} control, especially in P/D disaggregated serving. Our work focuses on deeply understanding the fine-grained frequency and routing control for varying loads and SLOs, and showing the benefits of using them with P/D disaggregated LLM serving.}
EcoServe~\cite{li2025ecoserve} focuses on reducing the operational and embodied carbon emissions throughout the hardware lifecycle. It proposes reusing, allocating and provisioning GPUs and CPUs based on offline/online inference, workload demand, model-size etc. TAPAS~\cite{stojkovic2025tapas} focused on routing requests to specific instances within rows/aisles of GPU data‑center racks by exploiting available thermal and power slack to avoid high-power hotspots. Recently, Heron\cite{reddy2025ai} proposes placing GPUs closer to renewable energy sources and adjust workloads for improved efficiency. These works are complementary to our contribution and underscore the importance of sustainable AI. \aryan{These systems operate primarily at the infrastructure, cluster, or lifecycle level, whereas \name focuses on runtime operating-point selection.} 
We deeply study frequency-control under varying loads and \name is the first to explore energy-efficient LLM serving under P/D-disaggregated architecture systematically with SLO-aware feedback baked into the system. 


\subsection{P/D Disaggregation in LLM Inference} 
\label{subsec:pd-dist}

To better manage the different compute characteristics of the prefill and decode phases, recent work has proposed prefill-decode (P/D) disaggregation, which assigns the two phases to separate GPU nodes. Systems such as SplitWise~\cite{patel2024splitwise}, TetriInfer~\cite{hu2024inference}, Llumnix ~\cite{sun2024llumnix} and DistServe~\cite{zhong2024distserve} show that this separation improves goodput and TTFT and ITL SLO attainment rates by avoiding phase-level contention and enabling distinct execution policies. \aryan{Our key distinction is that we treat P/D disaggregation not only as a latency-isolation mechanism, but also as an energy-control surface that exposes new opportunities for phase-specific adaptation.} As such, major open-source LLM libraries such as vLLM~\cite{kwon2023efficient} and SGLang~\cite{zheng2024sglang} have also added runtime support for this architectural pattern. However, these efforts focus primarily on improving metrics like latency and throughput. \aryan{In particular, prior P/D systems do not characterize the phase-specific U-shaped energy–frequency behavior or the decode-side batch-boundary inefficiencies that directly motivate our \governor and \router designs.} \aryan{To our knowledge, \name is the first P/D-disaggregated LLM serving system, to jointly combine phase-specific per-engine-iteration frequency control, online-adaptive latency prediction, and architecture-aware routing for both energy and system efficiency.}\sout{To our knowledge, they have neither explored the energy implications of P/D disaggregation nor exploited the new opportunities that this architectural separation enables for energy optimization, such as phase-specific \cameraready{per-engine-iteration} frequency control and energy-aware P/D routing policies.} 
